% Generated by the blueprint tool (blueprint-pdf/1) from a blueprint snapshot; edit the
% Markdown sources, not this file.  Compile with `latexmk -pdf` or by running
% `pdflatex` twice (the second run settles the contents and the page references).
% Images, if any, are read from assets/ next to this file.
\documentclass[11pt]{article}
\usepackage[T1]{fontenc}
\usepackage[utf8]{inputenc}
\usepackage{lmodern}
\usepackage{textcomp}
\usepackage[a4paper,margin=2.5cm,headheight=14pt]{geometry}
\usepackage{amsmath,amssymb,amsthm}
\usepackage{amscd}  % KaTeX's CD diagrams; part of amsmath, so always there
\usepackage{mathtools}
\usepackage{bm}
\usepackage{graphicx}
\usepackage{xcolor}
\usepackage{array,booktabs,tabularx}
\usepackage{fancyvrb}
\usepackage[normalem]{ulem}
\usepackage{microtype}
% Optional, and only with their fonts: a small TeX installation has the
% packages without the fonts, which fails at the first formula.
\IfFileExists{mathrsfs.sty}{\IfFileExists{rsfs10.tfm}{\usepackage{mathrsfs}}{}}{}
\IfFileExists{stmaryrd.sty}{\IfFileExists{stmary10.tfm}{\usepackage{stmaryrd}}{}}{}
\IfFileExists{cancel.sty}{\usepackage{cancel}}{}
\usepackage[hyphens]{url}
\usepackage[pdfusetitle,bookmarksnumbered=false]{hyperref}
\usepackage{bookmark}
\definecolor{bplink}{HTML}{1F4E9C}
\definecolor{bpmuted}{HTML}{5F676E}
\hypersetup{colorlinks=true,linkcolor=bplink,urlcolor=bplink,citecolor=bplink}
\setlength{\parindent}{0pt}
\setlength{\parskip}{0.55em plus 0.2em minus 0.1em}
% Lists as tight as the page's.
\makeatletter
\def\@listi{\leftmargin\leftmargini\parsep 0pt\topsep 0.3em plus 0.1em\itemsep 0.1em plus 0.1em}
\let\@listI\@listi
\def\@listii{\leftmargin\leftmarginii\labelwidth\leftmarginii\advance\labelwidth-\labelsep\parsep 0pt\topsep 0.1em\itemsep 0.1em}
% An ordered list that does not start at 1.
\newcommand{\bpstartat}[1]{\setcounter{\@enumctr}{\numexpr#1-1\relax}}
\makeatother
\setcounter{tocdepth}{3}
\emergencystretch=3em
\allowdisplaybreaks

% A numbered statement, set as a paper sets one: "Theorem 2.1 (Title)."
% #1 is the lead ("Theorem~2.1"), #2 what follows it in roman.
\newenvironment{bpstatement}[2]{\par\addvspace{0.6\baselineskip}\noindent\textbf{#1}#2\textbf{.}\enspace\ignorespaces}
  {\par\addvspace{0.4\baselineskip}}
\newenvironment{bpstatementit}[2]{\begin{bpstatement}{#1}{#2}\itshape\ignorespaces}{\end{bpstatement}}
% A step: a relation with prose of its own, read under the object it starts from.
\newenvironment{bpstep}[1]{\par\addvspace{0.3\baselineskip}\begin{list}{}{\leftmargin=1.6em\rightmargin=0pt\topsep=0pt}\item\relax
  {\small\color{bpmuted}#1}\par\nopagebreak}{\end{list}\par}
% Notes about where an object is written out.
\newenvironment{bpnote}{\par\small\color{bpmuted}\itshape}{\par}
\newcommand{\bpomitted}[1]{\par{\small\color{bpmuted}\itshape #1\par}}
\newcommand{\bpbodyheading}[1]{\par\addvspace{0.4\baselineskip}\noindent\textbf{#1}\par\nopagebreak}
\newcommand{\bprule}{\par\medskip\noindent\hfil\rule{0.3\linewidth}{0.4pt}\hfil\par\medskip}
\newcommand{\bpboundary}[1]{\par\nopagebreak{\small\color{bpmuted}#1\par}\nopagebreak}
% A reference to something this export leaves out, and one that resolves nowhere.
\newcommand{\bpout}[1]{#1\textsuperscript{\hyperlink{bp-outside}{\ensuremath{\ast}}}}
\newcommand{\bpbroken}[1]{\textcolor{red!70!black}{[#1?]}}
\newcommand{\bpslug}[1]{\texttt{#1}}
% A bibliography entry: label, anchor, entry; the label hangs in the margin.
\newcommand{\bpbibitem}[3]{\par\setbox0\hbox{[#1]\enspace}\ifdim\wd0<5em\setbox0\hbox to 5em{[#1]\hfil}\fi
  \hangindent=\wd0 \hangafter=1 \noindent\box0\phantomsection\label{#2}#3\par\smallskip}
% A character the export knows no way to set.
\newcommand{\bpglyph}[1]{\fbox{\scriptsize #1}}
% Images, at their own size unless they are wider than the line.
\newsavebox{\bpimagebox}
\newcommand{\bpimage}[1]{\sbox{\bpimagebox}{\includegraphics{#1}}%
  \ifdim\wd\bpimagebox>\linewidth\resizebox{\linewidth}{!}{\usebox{\bpimagebox}}\else\usebox{\bpimagebox}\fi}
\newcommand{\bpmissingimage}[1]{\fbox{\parbox{0.8\linewidth}{\centering\small\color{bpmuted}Image not included: #1}}}
\DefineVerbatimEnvironment{bpverbatim}{Verbatim}{fontsize=\small,xleftmargin=1em}
\renewcommand{\qedsymbol}{\ensuremath{\blacksquare}}
\hypersetup{pdftitle={Weil Conjectures — 7 The six operations on Dctf},pdfsubject={Chapter 7 The six operations on Dctf ; all levels},pdfcreator={blueprint}}
% KaTeX commands the website knows
\providecommand{\R}{\mathbb{R}}\providecommand{\N}{\mathbb{N}}\providecommand{\Z}{\mathbb{Z}}
\providecommand{\Q}{\mathbb{Q}}\providecommand{\C}{\mathbb{C}}
\providecommand{\Reals}{\mathbb{R}}\providecommand{\reals}{\mathbb{R}}\providecommand{\Complex}{\mathbb{C}}
\providecommand{\cnums}{\mathbb{C}}\providecommand{\natnums}{\mathbb{N}}\providecommand{\Bbb}[1]{\mathbb{#1}}
\providecommand{\bold}[1]{\mathbf{#1}}\providecommand{\frak}[1]{\mathfrak{#1}}\providecommand{\mathbbm}[1]{\mathbb{#1}}
\providecommand{\lang}{\langle}\providecommand{\rang}{\rangle}\providecommand{\bra}[1]{\langle #1\rvert}
\providecommand{\ket}[1]{\lvert #1\rangle}\providecommand{\braket}[1]{\langle #1\rangle}
\providecommand{\set}[1]{\{#1\}}\providecommand{\Set}[1]{\left\{#1\right\}}
\providecommand{\infin}{\infty}\providecommand{\isin}{\in}\providecommand{\plusmn}{\pm}\providecommand{\sub}{\subset}
\providecommand{\sube}{\subseteq}\providecommand{\supe}{\supseteq}\providecommand{\exist}{\exists}
\providecommand{\bull}{\bullet}\providecommand{\larr}{\leftarrow}\providecommand{\rarr}{\rightarrow}
\providecommand{\lrarr}{\leftrightarrow}\providecommand{\Larr}{\Leftarrow}\providecommand{\Rarr}{\Rightarrow}
\providecommand{\Lrarr}{\Leftrightarrow}\providecommand{\uarr}{\uparrow}\providecommand{\darr}{\downarrow}
\providecommand{\harr}{\leftrightarrow}\providecommand{\lArr}{\Leftarrow}\providecommand{\rArr}{\Rightarrow}
\providecommand{\hArr}{\Leftrightarrow}\providecommand{\alef}{\aleph}\providecommand{\alefsym}{\aleph}
\providecommand{\real}{\Re}\providecommand{\image}{\Im}\providecommand{\weierp}{\wp}
\providecommand{\Alpha}{\mathrm{A}}\providecommand{\Beta}{\mathrm{B}}\providecommand{\Epsilon}{\mathrm{E}}
\providecommand{\Zeta}{\mathrm{Z}}\providecommand{\Eta}{\mathrm{H}}\providecommand{\Iota}{\mathrm{I}}
\providecommand{\Kappa}{\mathrm{K}}\providecommand{\Mu}{\mathrm{M}}\providecommand{\Nu}{\mathrm{N}}
\providecommand{\Omicron}{\mathrm{O}}\providecommand{\Rho}{\mathrm{P}}\providecommand{\Tau}{\mathrm{T}}
\providecommand{\Chi}{\mathrm{X}}\providecommand{\omicron}{o}\providecommand{\gt}{>}\providecommand{\lt}{<}
\providecommand{\And}{\&}\providecommand{\hdashline}{\hline}\providecommand{\operatornamewithlimits}[1]{\operatorname*{#1}}
\providecommand{\llbracket}{[\![}\providecommand{\rrbracket}{]\!]}\providecommand{\cancel}[1]{#1}
\providecommand{\mathscr}{\mathcal}
% the characters beyond ASCII this document uses
\DeclareUnicodeCharacter{2160}{\textup{I}}
\DeclareUnicodeCharacter{2161}{\textup{II}}
% the project's macros ([katex.macros] in blueprint.toml)
\providecommand{\A}{}\renewcommand{\A}{\mathbf A}
\providecommand{\AEn}{}\renewcommand{\AEn}{A_{E,n}}
\providecommand{\Aaf}{}\renewcommand{\Aaf}{\mathbf A}
\providecommand{\Ab}{}\renewcommand{\Ab}{\mathsf{Ab}}
\providecommand{\AffSch}{}\renewcommand{\AffSch}{\operatorname{AffSch}}
\providecommand{\Aone}{}\renewcommand{\Aone}{\mathbf A^1}
\providecommand{\Bl}{}\renewcommand{\Bl}{\operatorname{Bl}}
\providecommand{\Br}{}\renewcommand{\Br}{\mathrm{Br}}
\providecommand{\C}{}\renewcommand{\C}{\mathbb C}
\providecommand{\Cech}{}\renewcommand{\Cech}{\check{H}}
\providecommand{\Coext}{}\renewcommand{\Coext}{\operatorname{Coext}}
\providecommand{\Cone}{}\renewcommand{\Cone}{\operatorname{Cone}}
\providecommand{\Cons}{}\renewcommand{\Cons}{\operatorname{Cons}}
\providecommand{\Cref}{}\renewcommand{\Cref}[1]{\text{#1}}
\providecommand{\D}{}\renewcommand{\D}{D}
\providecommand{\DLAFT}{}\renewcommand{\DLAFT}{\operatorname{FT}}
\providecommand{\DLAFr}{}\renewcommand{\DLAFr}{\operatorname{Fr}}
\providecommand{\DLAQ}{}\renewcommand{\DLAQ}{{\overline{\mathbf Q}_\ell}}
\providecommand{\DLRQ}{}\renewcommand{\DLRQ}{{\overline{\mathbf Q}_\ell}}
\providecommand{\Dbc}{}\renewcommand{\Dbc}{D^b_c}
\providecommand{\Dc}{}\renewcommand{\Dc}{D_{\mathrm{cons}}}
\providecommand{\Dctf}{}\renewcommand{\Dctf}{D^b_{ctf}}
\providecommand{\Dic}{}\renewcommand{\Dic}{D_{\IndCons}}
\providecommand{\Dplus}{}\renewcommand{\Dplus}{D^{+}}
\providecommand{\End}{}\renewcommand{\End}{\operatorname{End}}
\providecommand{\Et}{}\renewcommand{\Et}{\mathrm{\acute{e}t}}
\providecommand{\Ex}{}\renewcommand{\Ex}{\mathrm{Ex}}
\providecommand{\Ext}{}\renewcommand{\Ext}{\operatorname{Ext}}
\providecommand{\F}{}\renewcommand{\F}{\mathbb F}
\providecommand{\FT}{}\renewcommand{\FT}{\operatorname{FT}}
\providecommand{\Fbar}{}\renewcommand{\Fbar}{{\overline{\mathbf F}_q}}
\providecommand{\Fell}{}\renewcommand{\Fell}{\mathbf F_\ell}
\providecommand{\Fib}{}\renewcommand{\Fib}{\operatorname{Fib}}
\providecommand{\Fix}{}\renewcommand{\Fix}{\operatorname{Fix}}
\providecommand{\Fp}{}\renewcommand{\Fp}{{\mathbf F_{p}}}
\providecommand{\Fq}{}\renewcommand{\Fq}{{\mathbf F_q}}
\providecommand{\Fqn}{}\renewcommand{\Fqn}{{\mathbf F_{q^{n}}}}
\providecommand{\Frob}{}\renewcommand{\Frob}{\operatorname{Frob}}
\providecommand{\GL}{}\renewcommand{\GL}{\operatorname{GL}}
\providecommand{\Gal}{}\renewcommand{\Gal}{\operatorname{Gal}}
\providecommand{\Hens}{}\renewcommand{\Hens}{\operatorname{Hens}}
\providecommand{\Hilb}{}\renewcommand{\Hilb}{\operatorname{Hilb}}
\providecommand{\Hom}{}\renewcommand{\Hom}{\operatorname{Hom}}
\providecommand{\IndCons}{}\renewcommand{\IndCons}{\operatorname{IndCons}}
\providecommand{\Kar}{}\renewcommand{\Kar}{\operatorname{Kar}}
\providecommand{\LT}{}\renewcommand{\LT}{\operatorname{LT}}
\providecommand{\Lie}{}\renewcommand{\Lie}{\operatorname{Lie}}
\providecommand{\Loc}{}\renewcommand{\Loc}{\operatorname{Loc}}
\providecommand{\Mod}{}\renewcommand{\Mod}{\mathsf{Mod}}
\providecommand{\N}{}\renewcommand{\N}{\mathbb{N}}
\providecommand{\Needspace}{}\renewcommand{\Needspace}{\par}
\providecommand{\O}{}\renewcommand{\O}{{\mathcal{O}}}
\providecommand{\Ob}{}\renewcommand{\Ob}{\mathrm{Ob}}
\providecommand{\Oe}{}\renewcommand{\Oe}{\mathcal O_E}
\providecommand{\PP}{}\renewcommand{\PP}{\mathbf{P}}
\providecommand{\Perf}{}\renewcommand{\Perf}{\operatorname{Perf}}
\providecommand{\Pic}{}\renewcommand{\Pic}{\mathrm{Pic}}
\providecommand{\Pm}{}\renewcommand{\Pm}{\mathbf P^{m}}
\providecommand{\Pone}{}\renewcommand{\Pone}{\mathbf P^{1}}
\providecommand{\Pqn}{}\renewcommand{\Pqn}{\mathbf P^{n}}
\providecommand{\ProFin}{}\renewcommand{\ProFin}{\operatorname{ProFin}}
\providecommand{\Proet}{}\renewcommand{\Proet}{\mathrm{pro\acute{e}t}}
\providecommand{\Q}{}\renewcommand{\Q}{\mathbb Q}
\providecommand{\Qb}{}\renewcommand{\Qb}{\overline{\mathbf Q}_{\ell}}
\providecommand{\Qbar}{}\renewcommand{\Qbar}{\overline{\mathbf Q}}
\providecommand{\QbarX}{}\renewcommand{\QbarX}{\overline{\mathbf Q}_{\ell,X}}
\providecommand{\Qbarl}{}\renewcommand{\Qbarl}{{\overline{\mathbf Q}_{\ell}}}
\providecommand{\Qell}{}\renewcommand{\Qell}{{\mathbf{Q}_{\ell}}}
\providecommand{\Ql}{}\renewcommand{\Ql}{\mathbf Q_{\ell}}
\providecommand{\Qlb}{}\renewcommand{\Qlb}{{\overline{\mathbf Q}_\ell}}
\providecommand{\Qlbar}{}\renewcommand{\Qlbar}{\overline{\mathbf Q}_\ell}
\providecommand{\R}{}\renewcommand{\R}{\mathbb{R}}
\providecommand{\RG}{}\renewcommand{\RG}{\operatorname{\mathrm{R}\Gamma}}
\providecommand{\RGamma}{}\renewcommand{\RGamma}{R\Gamma}
\providecommand{\RGammac}{}\renewcommand{\RGammac}{R\Gamma_c}
\providecommand{\RHom}{}\renewcommand{\RHom}{R\mathcal{H}om}
\providecommand{\RHomX}{}\renewcommand{\RHomX}{R\mathcal H om_X}
\providecommand{\Rat}{}\renewcommand{\Rat}{\operatorname{IsRat}}
\providecommand{\Res}{}\renewcommand{\Res}{\operatorname{Res}}
\providecommand{\Rlim}{}\renewcommand{\Rlim}{R\!\varprojlim}
\providecommand{\SIm}{}\renewcommand{\SIm}{\operatorname{SIm}}
\providecommand{\Sch}{}\renewcommand{\Sch}{\mathsf{Sch}}
\providecommand{\Sh}{}\renewcommand{\Sh}{\operatorname{Sh}}
\providecommand{\Span}{}\renewcommand{\Span}{\operatorname{span}}
\providecommand{\Spec}{}\renewcommand{\Spec}{\operatorname{Spec}}
\providecommand{\Tr}{}\renewcommand{\Tr}{\operatorname{Tr}}
\providecommand{\Trsm}{}\renewcommand{\Trsm}{\operatorname{Tr}^{\mathrm{sm}}}
\providecommand{\Wt}{}\renewcommand{\Wt}{\operatorname{Wt}}
\providecommand{\Z}{}\renewcommand{\Z}{\mathbb{Z}}
\providecommand{\Zar}{}\renewcommand{\Zar}{\mathrm{zar}}
\providecommand{\Zell}{}\renewcommand{\Zell}{{\mathbf{Z}_{\ell}}}
\providecommand{\Zl}{}\renewcommand{\Zl}{{\mathbf Z_\ell}}
\providecommand{\affet}{}\renewcommand{\affet}[1]{{#1}^{\mathrm{aff}}_{\mathrm{ét}}}
\providecommand{\affproet}{}\renewcommand{\affproet}[1]{{#1}^{\mathrm{aff}}_{\mathrm{proét}}}
\providecommand{\boxt}{}\renewcommand{\boxt}{\boxtimes_A^{\mathrm L}}
\providecommand{\cF}{}\renewcommand{\cF}{\mathcal{F}}
\providecommand{\cHom}{}\renewcommand{\cHom}{R\mathcal{H}\!om}
\providecommand{\cI}{}\renewcommand{\cI}{\mathcal{I}}
\providecommand{\cL}{}\renewcommand{\cL}{\mathcal L}
\providecommand{\cM}{}\renewcommand{\cM}{\mathcal M}
\providecommand{\cN}{}\renewcommand{\cN}{\mathcal N}
\providecommand{\cO}{}\renewcommand{\cO}{\mathcal O}
\providecommand{\cT}{}\renewcommand{\cT}{\mathcal T}
\providecommand{\checked}{}\renewcommand{\checked}{\color{red}\aftergroup\restorecolor}
\providecommand{\coeff}{}\renewcommand{\coeff}{\operatorname{coeff}}
\providecommand{\coker}{}\renewcommand{\coker}{\operatorname{coker}}
\providecommand{\colim}{}\renewcommand{\colim}{\operatorname{colim}}
\providecommand{\constructedby}{}\renewcommand{\constructedby}[1]{\par\smallskip\noindent\textcolor{purple!60!black}{\textsf{Constructed by:} }
  \footnotesize\sloppy\nolinkurl{#1}\par}
\providecommand{\deepresult}{}\renewcommand{\deepresult}[1]{\par\smallskip\noindent\textcolor{orange!45!black}{
  \textsf{Blueprint role: named deep theorem node; see the cluster proof architecture.}}
  \par\noindent\textcolor{orange!45!black}{\textsf{Mathematical source:} }#1\par}
\providecommand{\dt}{}\renewcommand{\dt}{\mathbin{\otimes^{\mathbf L}}}
\providecommand{\dw}{}\renewcommand{\dw}{\operatorname{dw}}
\providecommand{\et}{}\renewcommand{\et}[1]{{#1}_{\mathrm{ét}}}
\providecommand{\ev}{}\renewcommand{\ev}{\operatorname{ev}}
\providecommand{\fib}{}\renewcommand{\fib}{\operatorname{fib}}
\providecommand{\formalizationnote}{}\renewcommand{\formalizationnote}[1]{\begin{quote}\small\raggedright\textbf{Formalization note.} #1\end{quote}}
\providecommand{\fpqc}{}\renewcommand{\fpqc}{X_{\mathrm{fpqc}}}
\providecommand{\gammacont}{}\renewcommand{\gammacont}{\Gamma_{\mathrm{cont}}}
\providecommand{\gm}{}\renewcommand{\gm}{\mathbf G_m}
\providecommand{\hocolim}{}\renewcommand{\hocolim}{\operatorname{hocolim}}
\providecommand{\holim}{}\renewcommand{\holim}{\operatorname{holim}}
\providecommand{\hrB}{}\renewcommand{\hrB}{\@ifnextchar]{\hrEnd}{\hrStep}}
\providecommand{\hrEnd}{}\renewcommand{\hrEnd}{\@firstoftwo{}}
\providecommand{\hrStep}{}\renewcommand{\hrStep}{\@firstoftwo{\hrB}}
\providecommand{\id}{}\renewcommand{\id}{\operatorname{id}}
\providecommand{\kunFl}{}\renewcommand{\kunFl}{\mathbf F_\ell}
\providecommand{\kunZn}{}\renewcommand{\kunZn}{\mathbf Z/n\mathbf Z}
\providecommand{\kuncd}{}\renewcommand{\kuncd}{\operatorname{cd}}
\providecommand{\kundim}{}\renewcommand{\kundim}{\dim}
\providecommand{\kunotimes}{}\renewcommand{\kunotimes}[1]{\otimes^{\mathbf L}_{#1}}
\providecommand{\kunsupp}{}\renewcommand{\kunsupp}{\operatorname{Supp}}
\providecommand{\lotimes}{}\renewcommand{\lotimes}{\otimes^{\mathrm L}}
\providecommand{\m}{}\renewcommand{\m}{{\mathfrak{m}}}
\providecommand{\n}{}\renewcommand{\n}{{\mathfrak{n}}}
\providecommand{\nagataGap}{}\renewcommand{\nagataGap}[1]{\par\noindent\textbf{Formalization gap.} #1\par}
\providecommand{\openobligation}{}\renewcommand{\openobligation}[1]{\par\smallskip\noindent\textcolor{red!65!black}{\textsf{Open obligation:} }
  \footnotesize\sloppy\nolinkurl{#1}\par}
\providecommand{\p}{}\renewcommand{\p}{{\mathfrak{p}}}
\providecommand{\pr}{}\renewcommand{\pr}{\operatorname{pr}}
\providecommand{\preshv}{}\renewcommand{\preshv}[1]{\mathrm{PreShv}(#1)}
\providecommand{\pro}{}\renewcommand{\pro}[1]{\mathrm{Pro}(#1)}
\providecommand{\proet}{}\renewcommand{\proet}{X_{\mathrm{proét}}}
\providecommand{\q}{}\renewcommand{\q}{{\mathfrak{q}}}
\providecommand{\rad}{}\renewcommand{\rad}{\operatorname{rad}}
\providecommand{\rank}{}\renewcommand{\rank}{\operatorname{rank}}
\providecommand{\red}{}\renewcommand{\red}{\color{red}}
\providecommand{\restorecolor}{}\renewcommand{\restorecolor}{\color{black}}
\providecommand{\retiredobligation}{}\renewcommand{\retiredobligation}[1]{\par\smallskip\noindent\textcolor{red!75!black}{\textsf{Retired invalid route:} }
  \footnotesize\sloppy\nolinkurl{#1}\par}
\providecommand{\rk}{}\renewcommand{\rk}{\operatorname{rk}}
\providecommand{\sep}{}\renewcommand{\sep}{\textnormal{sep}}
\providecommand{\shv}{}\renewcommand{\shv}[1]{\mathrm{Shv}(#1)}
\providecommand{\spec}{}\renewcommand{\spec}[1]{\mathrm{Spec}(#1)}
\providecommand{\stacks}{}\renewcommand{\stacks}[1]{\href{https://stacks.math.columbia.edu/tag/#1}{\textnormal{Stacks, Tag~#1}}}
\providecommand{\stacksproject}{}\renewcommand{\stacksproject}[1]{{\cite[\href{https://stacks.math.columbia.edu/tag/#1}{Tag #1}]{stacks-project}}}
\providecommand{\statusnote}{}\renewcommand{\statusnote}[1]{\begin{quote}\small\raggedright\textbf{Status.} #1\end{quote}}
\providecommand{\tildeideal}{}\renewcommand{\tildeideal}[2]{\widetilde{(#1, #2)}}
\providecommand{\tildering}{}\renewcommand{\tildering}[3]{{#1}_{\widetilde{(#2, #3)}}}
\providecommand{\tr}{}\renewcommand{\tr}{\operatorname{tr}}
\providecommand{\wt}{}\renewcommand{\wt}{\operatorname{wt}}

\begin{document}
\pagestyle{plain}
\begin{center}
{\LARGE\bfseries Weil Conjectures\par}
\vspace{0.7em}
{\large Chapter 7: The six operations on $\Dctf$\par}
\vspace{0.5em}
{\small\color{bpmuted}Snapshot 2026-09-26, revision 8e4cb88e2a\par}
\end{center}
\vspace{0.5em}
{\small\tableofcontents}
\bigskip\hrule\bigskip

% object: chap-The-six-operations-on-Dbctf (blueprint/chap-The-six-operations-on-Dbctf/_section.md)
\section*{\texorpdfstring{7\quad The six operations on $\Dctf$}{7 The six operations on Dctf}}
\phantomsection\label{bp1}%
\addcontentsline{toc}{section}{\texorpdfstring{7\quad The six operations on $\Dctf$}{7 The six operations on Dctf}}
% text: prose-07-01
\phantomsection\label{bp2}%
Smooth base change and Verdier biduality will be proved with torsion coefficients, as were proper base change and Deligne's theorem on \(Rj_*\). The Weil conjectures, however, need them over \(\overline{\mathbf Q}_\ell\), where Frobenius eigenvalues live. This chapter provides the principle that carries every such statement across. A canonical map built from the six operations is an isomorphism for all admissible \(\Lambda\) as soon as it is one for every finite ring (Proposition \hyperref[bp34]{7.12}).

For \(\mathcal O_E\) both sides are derived complete, and by derived Nakayama it suffices to check the map modulo \(\pi_E\). For instance, let \(K\) be in \(D^b_{\mathrm{ctf}}(X,\mathcal O_E)\). Both sides of the biduality map \(K\to\mathbb D_X\mathbb D_XK\) are derived complete. Modulo \(\pi_E\), moreover, the map is biduality with \(\mathbf F_E\)-coefficients, which is an isomorphism. Hence the biduality map is an isomorphism. From there we pass to \(E\) and \(\overline{\mathbf Q}_\ell\): every object over \(E\) is built by cones from integral models with \(\ell\) inverted, and every object over \(\overline{\mathbf Q}_\ell\) is defined over a finite extension \(E\) (Proposition \hyperref[bp29]{7.11}). In this way smooth base change (Theorem \bpout{8.11}) and Verdier biduality (Theorem \bpout{10.8}) pass from torsion to \(\ell\)-adic coefficients.

The principle needs two things, and neither is formal. First, both sides must lie in \(D^b_{\mathrm{ctf}}(-,\Lambda)\). Thus all six operations must preserve these categories, and not only \(Rf_!\) and \(f^!\), which we now control (Theorem \bpout{5.2}, Section \bpout{6.8}). Second, each operation must commute with change of coefficients. For \(f^*\) and \(\otimes^{\mathbf L}\) this is clear, since a change of coefficients is itself a tensor product. For a right adjoint, however, it is a real statement. For example, consider a map of finite rings \(A\to B\); then the isomorphism \(R\Gamma(X,K)\otimes^{\mathbf L}_AB\simeq
R\Gamma(X,K\otimes^{\mathbf L}_AB)\) is a projection formula, and it uses finite cohomological dimension. Next consider \(\mathcal O_E\to E\). Since \(E\) is the colimit of the modules \(\pi_E^{-n}\mathcal O_E\), the comparison says that \(Rf_*\) commutes with a filtered colimit. This holds because the terms are uniformly bounded below.

Internal Hom causes a further difficulty, because it hides the exceptional inverse image. For a closed immersion \(i\) we have

\[
\cHom(i_*\Lambda,K)\simeq i_*i^!K .
\]

Before we can say anything about \(\cHom(K,L)\otimes^{\mathbf L}\Lambda'\), we must therefore know how \(i^!\) behaves under change of coefficients. The comparison map for \(f^!\) is defined by adjunction, and to show that it is an isomorphism we compute \(f^!\) once more, for a closed immersion and for \(\mathbf P^1\). For \(\mathbf P^1\), for instance, it comes down to the compatibility of the cycle class of the diagonal and of the trace with change of coefficients (Proposition \hyperref[bp26]{7.10}).

The first requirement needs no construction: on the pro-étale site the six operations are the ordinary derived functors of sheaves of modules, and we only check that they preserve \(D^b_{\mathrm{ctf}}(-,\Lambda)\) (Proposition \hyperref[bp5]{7.1}). For torsion coefficients on the étale site, these finiteness theorems are in SGA 4 and SGA 4\(\tfrac12\) [\hyperref[bp37]{SGA4}], [\hyperref[bp38]{SGA4½}]. In general, the work is in \(Rf_*\) and in \(\cHom\) (Section \hyperref[bp3]{7.1}). There, a compactification reduces \(Rf_*\) to proper direct image and \(Rj_*\), while the formula above and its analogue for open immersions express \(\cHom\) through \(Ra_*\), \(a^!\) and tensor products. Tensor products, finally, are handled by the local description of a complex of finite Tor amplitude: on the strata of a finite stratification it is locally \(\Lambda_S\otimes P\) with \(P\) perfect (Section \bpout{3.4}).

From now on we work in \(D^b_{\mathrm{ctf}}\) without further comment and pass between \(\mathcal O_E\), \(E\), \(\overline{\mathbf Q}_\ell\) and finite rings whenever a statement is easier for one of them, as in the trace formula and in the proof of the Weil conjectures.

% object: sec-Stability-of-Dbctf (blueprint/chap-The-six-operations-on-Dbctf/sec-Stability-of-Dbctf/_section.md), content hidden
\subsection*{\texorpdfstring{7.1\quad Stability of $\Dctf$}{7.1 Stability of Dctf}}
\phantomsection\label{bp3}%
\addcontentsline{toc}{subsection}{\texorpdfstring{7.1\quad Stability of $\Dctf$}{7.1 Stability of Dctf}}
% text: prose-07-02
\phantomsection\label{bp4}%
Pullback and tensor product are local, and they preserve the description of a complex of \(D^b_{\mathrm{ctf}}\) on the strata of a stratification as \(\Lambda_S\otimes P\) with \(P\) perfect. The direct images and internal Hom are not local. For the direct images a compactification \(f=pj\) helps. Since \(Rf_*=Rp_*Rj_*\), it suffices to know proper direct image and direct image along an open immersion. Internal Hom hides both of these operations and also \(f^!\). For an open immersion \(j\) we have \(\cHom(j_!\Lambda,K)=Rj_*j^*K\), and for a closed immersion \(i\) we have the formula \(\cHom(i_*\Lambda,K)\simeq i_*i^!K\) recalled at the beginning of this chapter. If \(a\) is the locally closed immersion of a stratum and \(P\) is a perfect complex on it, the formula

\[
\cHom(a_!P,K)\simeq Ra_*\bigl(P^\vee\otimes^{\mathbf L}a^!K\bigr)
\]

combines the two. Every constructible complex is built from objects \(a_!P\) by finitely many cones. By this formula, internal Hom is therefore expressed through \(Ra_*\), \(a^!\) and tensor products. This is one reason we needed \(f^!\) on \(D^b_{\mathrm{ctf}}\) first. The result is Proposition \hyperref[bp5]{7.1}.

% object: proposition-six-operations-dctf (blueprint/chap-The-six-operations-on-Dbctf/sec-Stability-of-Dbctf/proposition-six-operations-dctf/_section.md)
\phantomsection\label{bp5}%
\begin{bpstatementit}{Proposition~7.1}{ (The six operations on $\Dctf$)}
For every admissible \(\Lambda\in\{A,\mathbf Z_\ell,\mathcal O_E,\mathbf Q_\ell,E,\overline{\mathbf Q}_\ell\}\), the operations \(f^*\), \(Rf_*\), \(Rf_!\), \(f^!\), \(\otimes^{\mathbf L}\) and \(\cHom\) preserve \(D^b_{\mathrm{ctf}}(-,\Lambda)\) for morphisms of schemes of finite type over \(k\). Their adjunctions are the restrictions of the adjunctions \(f^*\dashv Rf_*\) and \(-\otimes^{\mathbf L}K\dashv\cHom(K,-)\) on \(D(-_{\mathrm{pro}\acute{e}t},\Lambda)\) or of the adjunction \(Rf_!\dashv f^!\) on \(\Loc(D^b_{\mathrm{ctf}}(-,\Lambda))\).
\end{bpstatementit}

% proof: proof/proposition-six-operations-dctf
\begin{proof}[{Proof}]\phantomsection\label{bp6}%
We treat the operations in the order in which they depend on each other: first the local operations \(f^*\) and \(\otimes^{\mathbf L}\), then the two direct images, and finally internal Hom, which we express through the others.

Pullback and tensor product preserve the property of being, on the strata of a finite stratification, locally \(\Lambda_S\otimes P\) with \(P\) a perfect complex. For pullback this is clear, since the strata pull back to a stratification, while for tensor products we take a common refinement of the two stratifications, and on each of its strata we use that \(P\otimes^{\mathbf L}Q\) is perfect when \(P\) and \(Q\) are. Similarly, extension by zero preserves finite Tor amplitude, since it is exact and satisfies the projection formula \(a_!K\otimes^{\mathbf L}N\simeq a_!(K\otimes^{\mathbf L}a^*N)\). Concretely, for finite \(A\) these statements can be checked on the bounded flat constructible models of Lemma \bpout{3.13}. For \(\mathcal O_E\), \(E\) and \(\overline{\mathbf Q}_\ell\), on the other hand, they follow from Lemma \bpout{5.20} and from Definition \bpout{3.1} of constructible sheaves: on each stratum, such a sheaf is locally \(\Lambda_X\otimes_{\Lambda(*)}P\) with \(P\) finitely generated over the discrete valuation ring or field \(\Lambda(*)\). Over such a ring a finitely generated module has a finite free resolution of length at most one. Hence \(P\) is perfect.

For the direct images we use a compactification to separate the proper part from the open part. Proper direct image preserves \(D^b_{\mathrm{ctf}}(-,\Lambda)\) by Proposition \bpout{4.40}. Direct image along an open immersion preserves it by Proposition \bpout{6.19} for finite \(A\) and Proposition \bpout{6.21} for \(\mathcal O_E\), \(E\) and \(\overline{\mathbf Q}_\ell\). If \(f=pj\) is a compactification, then \(Rf_*=Rp_*Rj_*\) and \(Rf_!=Rp_*j_!\). Each of these composites consists of functors already known to preserve \(D^b_{\mathrm{ctf}}\), which proves both assertions on direct images. That \(f^!\) preserves \(D^b_{\mathrm{ctf}}(-,\Lambda)\) is part of Theorem \bpout{5.2}.

For internal Hom we reduce to a first argument supported on a single stratum, where \(\cHom\) can be written in terms of \(Ra_*\), \(a^!\) and \(\otimes^{\mathbf L}\). This is possible because the functor \(\cHom(-,K)\) sends triangles in the first argument to triangles, and \(D^b_{\mathrm{ctf}}\) is closed under cones and shifts. Hence, for a triangle \(M'\to M\to M''\), the complex \(\cHom(M,K)\) lies in \(D^b_{\mathrm{ctf}}\) whenever \(\cHom(M',K)\) and \(\cHom(M'',K)\) do. Using a finite stratification and finite cones, we may therefore reduce to the case where the first argument is \(a_!P\). Here \(a:S\hookrightarrow X\) is locally closed and \(P\) is a locally constant perfect complex on \(S\). For finite \(A\) this reduction follows from the bounded flat model and its stupid filtration, and for \(\mathcal O_E\), \(E\) and \(\overline{\mathbf Q}_\ell\) it follows from the same description of constructible sheaves on strata. By tensor-Hom adjunction and the projection formula for \(a_!\) we have

\[
\cHom_X(a_!P,K)
\simeq Ra_*\cHom_S(P,a^!K)
\simeq Ra_*(P^\vee\otimes^{\mathbf L}a^!K).
\]

Here \(P^\vee=\cHom_S(P,\Lambda_S)\), and the second identity holds because \(P\) is dualizable. The last term belongs to \(D^b_{\mathrm{ctf}}(-,\Lambda)\) by the cases already established, namely those of \(a^!\), of tensor product and of \(Ra_*\).

Finally, since the inclusions are full, the Hom bijections of the ambient categories restrict to the stated adjunctions.
\end{proof}

% object: example-internal-hom-affine-line (blueprint/chap-The-six-operations-on-Dbctf/sec-Stability-of-Dbctf/example-internal-hom-affine-line/_section.md)
\phantomsection\label{bp7}%
\begin{bpstatement}{Example~7.2}{ (Internal Hom on the affine line)}
Over an algebraically closed \(k\), let \(j:U=\mathbf G_{m,k}\hookrightarrow
X=\mathbf A^1_k\) and let \(i\) be the inclusion of the origin. The localization sequence \(0\to j_!\Lambda_U\to\Lambda_X\to i_*\Lambda\to0\) shows that \(\Lambda_X\) is built from \(j_!\Lambda_U\) and \(i_*\Lambda\), and the formula above computes \(\cHom\) out of each piece. By Example \bpout{6.25},

\[
\cHom(j_!\Lambda_U,\Lambda_X)=Rj_*\Lambda_U,
\qquad
\cHom(i_*\Lambda,\Lambda_X)=i_*i^!\Lambda_X\simeq i_*\Lambda(-1)[-2].
\]

The first has \(\mathcal H^0=\Lambda_X\) and \(\mathcal H^1=i_*\Lambda(-1)\). In the second, the local Ext sheaves from the skyscraper sheaf \(i_*\Lambda\) to the constant sheaf vanish except in degree two. Both results are constructible but not lisse. The first is computed by a nonproper direct image, and the second by \(i^!\).
\end{bpstatement}


% object: sec-Change-of-coefficients (blueprint/chap-The-six-operations-on-Dbctf/sec-Change-of-coefficients/_section.md), content hidden
\subsection*{\texorpdfstring{7.2\quad Change of coefficients}{7.2 Change of coefficients}}
\phantomsection\label{bp8}%
\addcontentsline{toc}{subsection}{\texorpdfstring{7.2\quad Change of coefficients}{7.2 Change of coefficients}}
% text: prose-07-03
\phantomsection\label{bp9}%
Smooth base change for \(E\) and \(\overline{\mathbf Q}_\ell\) is deduced from the integral case by extending scalars (Chapter \bpout{8}), and later proofs often argue in the same way, through the reduction principle of Proposition \hyperref[bp34]{7.12}. For this the six operations must commute with extension of scalars. For pullback, extension by zero and tensor product this is formal, since extension of scalars is itself a tensor product. For \(Rf_*\), by contrast, it rests on finite cohomological dimension and on continuity, and for \(f^!\) it rests on a second computation of \(f^!\), for closed immersions and for \(\mathbf P^1\). These results are collected in Proposition \hyperref[bp26]{7.10}.

A change of coefficients is derived extension of scalars along a map of condensed rings, pulled back to \(X\). For a map \(\Lambda\to\Lambda'\) write

\[
K_{\Lambda'}=\Lambda'_X\otimes^{\mathbf L}_{\Lambda_X}K.
\]

This functor \(D(X_{\mathrm{pro}\acute{e}t},\Lambda)\to D(X_{\mathrm{pro}\acute{e}t},\Lambda')\) is left adjoint to the exact restriction of scalars \(D(X_{\mathrm{pro}\acute{e}t},\Lambda')\to D(X_{\mathrm{pro}\acute{e}t},\Lambda)\). We consider the following maps of coefficient rings:

\begin{enumerate}
\item any homomorphism \(A\to B\) of finite admissible rings;
\item the maps \(\mathcal O_E\to\mathcal O_F\) and \(E\to F\) for a finite extension \(F/E\) inside \(\overline{\mathbf Q}_\ell\);
\item localization \(\mathcal O_E\to E\), inclusions \(E\to\overline{\mathbf Q}_\ell\), and continuous maps \(\mathcal O_E\to A\) to finite admissible rings, such as the reductions \(\mathcal O_E\to\mathcal O_E/\pi_E^n\) (continuity is needed for the map to define a map of condensed rings);
\item composites of these maps.
\end{enumerate}

For \(E=\mathbf Q_\ell\), the list contains \(\mathbf Z_\ell\to\mathcal O_F\), \(\mathbf Q_\ell\to F\), \(\mathbf Z_\ell\to\mathbf Q_\ell\), \(\mathbf Z_\ell\to A\) and \(\mathbf Q_\ell\to\overline{\mathbf Q}_\ell\). Thus these maps connect all the admissible coefficient rings. Note that the map \(A\to B\) in (1) need not be flat, and for this reason extension of scalars is derived. As a consequence, the comparison isomorphisms below hold for complexes but not degree by degree, as the following example shows.

% object: example-derived-reduction-point (blueprint/chap-The-six-operations-on-Dbctf/sec-Change-of-coefficients/example-derived-reduction-point/_section.md)
\phantomsection\label{bp10}%
\begin{bpstatement}{Example~7.3}{ (Reducing cohomology modulo $\ell$)}
Let \(X=\operatorname{Spec}\mathbf F_q\) and consider \(\mathbf Z/\ell^2\to\mathbf F_\ell\). Let \(K\) be the lisse sheaf \(\mathbf Z/\ell^2\) on which Frobenius, the topological generator of \(\operatorname{Gal}(\overline{\mathbf F}_q/\mathbf F_q)
\simeq\widehat{\mathbf Z}\), acts by multiplication by the unit \(1+\ell\). Recall that for a finite module \(M\) with a continuous action of \(\widehat{\mathbf Z}\), cohomology is computed by the complex \(M\xrightarrow{F-1}M\) in degrees \(0\) and \(1\), where \(F\) is the action of the generator. For \(K\) this differential is multiplication by \(\ell\). Hence \(H^0(X,K)=\ell\mathbf Z/\ell^2\) and \(H^1(X,K)=(\mathbf Z/\ell^2)/\ell\), both of order \(\ell\). On the other hand, \(K_{\mathbf F_\ell}=K/\ell K\) is \(\mathbf F_\ell\) with trivial action, and \(H^0(X,K/\ell K)=\mathbf F_\ell\). The map \(H^0(X,K)\otimes\mathbf F_\ell\to H^0(X,K/\ell K)\) is zero, since the generator \(\ell\) of \(H^0(X,K)\) reduces to \(0\).

The derived comparison is nevertheless an isomorphism. The terms of the complex are free over \(\mathbf Z/\ell^2\). Therefore \(R\Gamma(X,K)\otimes^{\mathbf L}\mathbf F_\ell\) is \(\mathbf F_\ell\xrightarrow{0}\mathbf F_\ell\), which is \(R\Gamma(X,K/\ell K)\). The class \(1\in H^0(X,K/\ell K)\) comes from the element \(1\) of the degree-zero term, which is not a cocycle before reduction.
\end{bpstatement}

% text: prose-07-04
\phantomsection\label{bp11}%
Before stating the result we fix the comparison maps. Items (1)–(3) below, for \(f^*\), \(j_!\) and \(\otimes^{\mathbf L}\), are identifications of two composites of left adjoints. The remaining items are derived from these: items (4) and (6), for \(Rf_*\) and \(\cHom\), are adjoint transposes of the inverse of (1), respectively (3), composed with the extension of scalars of the counit, and item (5) combines (4) for \(p\) with (2).

% object: definition-change-of-coefficients-maps (blueprint/chap-The-six-operations-on-Dbctf/sec-Change-of-coefficients/definition-change-of-coefficients-maps/_section.md)
\phantomsection\label{bp12}%
\begin{bpstatement}{Definition~7.4}{ (Change-of-coefficients maps)}
Let \(\Lambda\to\Lambda'\) be one of the maps listed above, and let \(f:X\to Y\) be a morphism of schemes of finite type over \(k\).

\begin{enumerate}
\item For \(K\in D(Y_{\mathrm{pro}\acute{e}t},\Lambda)\), the isomorphism \((f^*K)_{\Lambda'}\simeq f^*(K_{\Lambda'})\) is the one given by associativity of derived pullback for ringed topoi [Stacks, \href{https://stacks.math.columbia.edu/tag/06YV}{06YV}]: both sides are the derived pullback along the morphism of ringed topoi \((X_{\mathrm{pro}\acute{e}t},\Lambda'_X)\to(Y_{\mathrm{pro}\acute{e}t},\Lambda_Y)\), factored in the two possible ways.
\item For an open immersion \(j:U\to X\) and \(K\in D(U_{\mathrm{pro}\acute{e}t},\Lambda)\), the isomorphism \((j_!K)_{\Lambda'}\simeq j_!(K_{\Lambda'})\) is the isomorphism of left adjoints corresponding to the equality of their right adjoints, restriction of scalars followed by \(j^*\) and \(j^*\) followed by restriction of scalars.
\item For \(K,L\in D(X_{\mathrm{pro}\acute{e}t},\Lambda)\), the isomorphism \((K\otimes^{\mathbf L}_\Lambda L)_{\Lambda'}\simeq
  K_{\Lambda'}\otimes^{\mathbf L}_{\Lambda'}L_{\Lambda'}\) is the associativity isomorphism: both sides are \(K\otimes^{\mathbf L}_\Lambda L\otimes^{\mathbf L}_\Lambda\Lambda'\).
\item For \(K\in D(X_{\mathrm{pro}\acute{e}t},\Lambda)\), the map \((Rf_*K)_{\Lambda'}\to Rf_*(K_{\Lambda'})\) is the adjoint transpose, under \(f^*\dashv Rf_*\) on \(D(-_{\mathrm{pro}\acute{e}t},\Lambda')\), of \[
   f^*\bigl((Rf_*K)_{\Lambda'}\bigr)\simeq(f^*Rf_*K)_{\Lambda'}
   \xrightarrow{(\varepsilon_K)_{\Lambda'}}K_{\Lambda'},
   \] where the first arrow is the inverse of (1) and \(\varepsilon\) is the counit of \(f^*\dashv Rf_*\) on \(D(-_{\mathrm{pro}\acute{e}t},\Lambda)\).
\item If \(f\) is separated, \(f=p\circ j\) is a compactification and \(K\in D^b_{\mathrm{ctf}}(X,\Lambda)\), the map \((Rf_!K)_{\Lambda'}\to Rf_!(K_{\Lambda'})\) is the composite \[
   (Rp_*j_!K)_{\Lambda'}\longrightarrow Rp_*\bigl((j_!K)_{\Lambda'}\bigr)
   \simeq Rp_*j_!(K_{\Lambda'})
   \] of (4) for \(p\) and \(Rp_*\) of (2).
\item For \(K,L\in D(X_{\mathrm{pro}\acute{e}t},\Lambda)\), the map \(\cHom(K,L)_{\Lambda'}\to\cHom(K_{\Lambda'},L_{\Lambda'})\) is the adjoint transpose, under \(-\otimes^{\mathbf L}K_{\Lambda'}\dashv
  \cHom(K_{\Lambda'},-)\), of \[
   \cHom(K,L)_{\Lambda'}\otimes^{\mathbf L}_{\Lambda'}K_{\Lambda'}
   \simeq\bigl(\cHom(K,L)\otimes^{\mathbf L}_\Lambda K\bigr)_{\Lambda'}
   \xrightarrow{(\mathrm{ev})_{\Lambda'}}L_{\Lambda'},
   \] where the first arrow is the inverse of (3) and \(\mathrm{ev}\) is evaluation.
\end{enumerate}
\end{bpstatement}

% text: prose-07-05
\phantomsection\label{bp13}%
The map for \(f^!\) is defined in the same way, as the transpose of a counit, once we know that (5) is an isomorphism; see (7.1) below. We now show, one operation at a time, that these maps are isomorphisms on \(D^b_{\mathrm{ctf}}\).

In the lemmas below it suffices to treat the maps in items (1)–(3) of the list of coefficient maps, that is, all maps except composites. Indeed, extension of scalars along a composite of two maps of the list is the composite of the two extensions. Each map in Definition \hyperref[bp12]{7.4} for the composite is the composite of the maps for the two factors. For (1)–(3) this holds by the coherence of the associativity isomorphisms, and for (4)–(6) it holds because the adjoint transpose of a composite of such maps is the composite of the transposes.

% object: lemma-extension-scalars-preserves-dctf (blueprint/chap-The-six-operations-on-Dbctf/sec-Change-of-coefficients/lemma-extension-scalars-preserves-dctf/_section.md)
\phantomsection\label{bp14}%
\begin{bpstatementit}{Lemma~7.5}{ (Extension of scalars preserves $\Dctf$)}
For each map \(\Lambda\to\Lambda'\) listed above, extension of scalars sends \(D^b_{\mathrm{ctf}}(X,\Lambda)\) into \(D^b_{\mathrm{ctf}}(X,\Lambda')\).
\end{bpstatementit}

% proof: proof/lemma-extension-scalars-preserves-dctf
\begin{proof}[{Proof}]\phantomsection\label{bp15}%
An object that is locally \(\Lambda_S\otimes P\) on a stratum \(S\), with \(P\) perfect, remains of this form after derived extension of scalars, since \(P\otimes^{\mathbf L}_\Lambda\Lambda'\) is a perfect \(\Lambda'\)-complex. Moreover, extension of scalars commutes with extension by zero and with finite cones. Every object of \(D^b_{\mathrm{ctf}}(X,\Lambda)\) is built from such objects by extension by zero from strata and finite cones, and so its extension of scalars lies in \(D^b_{\mathrm{ctf}}(X,\Lambda')\).
\end{proof}

% object: lemma-change-coefficients-direct-image (blueprint/chap-The-six-operations-on-Dbctf/sec-Change-of-coefficients/lemma-change-coefficients-direct-image/_section.md)
\phantomsection\label{bp16}%
\begin{bpstatementit}{Lemma~7.6}{ (Change of coefficients for direct image)}
For each map \(\Lambda\to\Lambda'\) listed above, every morphism \(f:X\to Y\) of schemes of finite type over \(k\) and every \(K\in D^b_{\mathrm{ctf}}(X,\Lambda)\), the map \((Rf_*K)_{\Lambda'}\to Rf_*(K_{\Lambda'})\) of Definition definition-change-of-coefficients-maps is an isomorphism.
\end{bpstatementit}

% proof: proof/lemma-change-coefficients-direct-image
\begin{proof}[{Proof}]\phantomsection\label{bp17}%
The argument depends on the kind of map \(\Lambda\to\Lambda'\). For maps of finite rings we use a projection formula, whereas in the other cases extension of scalars is either a finite construction (finite sums and a cone), with which every exact functor commutes, or a filtered colimit, with which \(Rf_*\) commutes on uniformly bounded below systems.

\emph{Maps \(A\to B\) of finite rings.} We compute on the étale site using Proposition \bpout{5.15}, which identifies the pro-étale and étale categories in this case, and we test the comparison on derived sections over quasi-compact étale objects. For a quasi-compact étale \(V\to Y\), we apply the projection formula for constant coefficients [Stacks, \href{https://stacks.math.columbia.edu/tag/0F12}{0F12}] on both \(V\) and \(X_V=X\times_YV\) and obtain

\[
R\Gamma(V,(Rf_*K)_B)
\simeq R\Gamma(X_V,K)\otimes_A^{\mathbf L}B
\simeq R\Gamma(X_V,K_B).
\]

Here the first identification uses the projection formula on \(V\) and then \(R\Gamma(V,Rf_*K)=R\Gamma(X_V,K)\), and the second is the projection formula on \(X_V\). This projection formula needs only cohomological bounds for \(\ell\)-primary coefficients, which hold for schemes of finite type over \(k\), and so it applies to every finite \(A\)-algebra \(B\). Since these isomorphisms commute with restriction in \(V\), the comparison map is an isomorphism in \(D(Y_{\acute{e}t},B)\), and hence, by Proposition \bpout{5.15}, in \(D(Y_{\mathrm{pro}\acute{e}t},B)\).

\emph{Maps \(\mathcal O_E\to A\) with \(A\) finite.} Since \(\mathcal O_E\) is a discrete valuation ring, \(A\) has a finite free \(\mathcal O_E\)-resolution of length at most one. Its condensed version is again exact. Indeed, a continuous map from a profinite set to the finite quotient \(A\) is constant on the pieces of a finite clopen partition, and hence lifts to a continuous map to the free module. The pullback of this resolution to \(X\) is a finite free resolution of \(A_X\) over \(\mathcal O_{E,X}\). Consequently, tensoring with \(A_X\) is computed by finite sums and a cone, and these commute with the exact functor \(Rf_*\).

\emph{Finite extensions.} For \(\mathcal O_E\to\mathcal O_F\) and \(E\to F\), the ring \(\mathcal O_F\) is free of rank \([F:E]\) over \(\mathcal O_E\). Hence \(K_{\mathcal O_F}\) is, as an \(\mathcal O_E\)-complex, a finite direct sum of copies of \(K\). The assertion follows because \(Rf_*\) commutes with finite sums.

\emph{The maps \(\mathcal O_E\to E\) and \(E\to\overline{\mathbf Q}_\ell\).} These are not finite, and we use continuity. For \(\mathcal O_E\to E\), the complex \(K_E\) is the colimit of \(K\xrightarrow{\ell}K\xrightarrow{\ell}\cdots\), and for \(E\to\overline{\mathbf Q}_\ell\), the complex \(K_{\overline{\mathbf Q}_\ell}\) is the filtered colimit of the \(K_F\) over finite extensions \(F/E\) in \(\overline{\mathbf Q}_\ell\). Since \(K\) is bounded, all terms of these systems are bounded below by a common integer. Hence the colimits commute with \(Rf_*\) by Corollary \bpout{2.29}.
\end{proof}

% object: lemma-change-coefficients-compact-support (blueprint/chap-The-six-operations-on-Dbctf/sec-Change-of-coefficients/lemma-change-coefficients-compact-support/_section.md)
\phantomsection\label{bp18}%
\begin{bpstatementit}{Lemma~7.7}{ (Change of coefficients for compactly supported direct image)}
For each map \(\Lambda\to\Lambda'\) listed above, every separated morphism \(f:X\to Y\) of finite type and every \(K\in D^b_{\mathrm{ctf}}(X,\Lambda)\), the map \((Rf_!K)_{\Lambda'}\to Rf_!(K_{\Lambda'})\) of Definition definition-change-of-coefficients-maps is an isomorphism, and it does not depend on the compactification.
\end{bpstatementit}

% proof: proof/lemma-change-coefficients-compact-support
\begin{proof}[{Proof}]\phantomsection\label{bp19}%
The map is the composite of the isomorphism of Lemma \hyperref[bp16]{7.6} for \(p\), applied to \(j_!K\in D^b_{\mathrm{ctf}}\) (Proposition \bpout{5.22}), and of \(Rp_*\) of the isomorphism of Definition definition-change-of-coefficients-maps. For independence, let \(h:Z'\to Z\) be a refinement of compactifications. Both comparison maps (4) and (2) commute with the refinement map (5.7). Indeed, they are formed from identifications of pullbacks and adjunction maps. A map out of an extension by zero is determined by its restriction to the open subscheme, and there all maps involved are identities. We also use that (4) is compatible with composition of morphisms. Under \(R(ph)_*\simeq Rp_*Rh_*\), the map (4) for \(ph\) is \(Rp_*\) of (4) for \(h\) followed by (4) for \(p\), since transposes of composites are composites. As in the proof of Proposition \bpout{5.26}, this gives independence for any two compactifications.
\end{proof}

% text: prose-07-06
\phantomsection\label{bp20}%
We can now define the comparison map for \(f^!\). Let \(f:X\to Y\) be separated of finite type and \(K\in D^b_{\mathrm{ctf}}(Y,\Lambda)\). Then \(f^!K\in D^b_{\mathrm{ctf}}(X,\Lambda)\) by Theorem \bpout{5.2}, and we let

\[
\beta_{f,K}:(f^!K)_{\Lambda'}\longrightarrow f^!(K_{\Lambda'})
\tag{7.1}
\]

be the transpose, under the adjunction \(Rf_!\dashv f^!\) on \(\Loc(D^b_{\mathrm{ctf}}(-,\Lambda'))\), of the morphism

\[
Rf_!((f^!K)_{\Lambda'})
\xrightarrow{\sim}(Rf_!f^!K)_{\Lambda'}
\xrightarrow{(\varepsilon_K)_{\Lambda'}}K_{\Lambda'},
\]

whose first arrow is the inverse of the isomorphism of Lemma \hyperref[bp18]{7.7}.

The proof below uses three properties of \(\beta\). First, the maps \(\beta\) are compatible with composition: under \((gf)^!\simeq f^!g^!\) we have \(\beta_{gf,K}=f^!(\beta_{g,K})\circ\beta_{f,g^!K}\). This holds because the counit of a composite adjunction is the composite of the counits, and because the comparison for \(Rf_!\) is compatible with the composition isomorphisms of Proposition \bpout{5.28}; both are determined by their restrictions to the open subscheme. Second, let \(\Lambda\to\Lambda'\to\Lambda''\) be a composite of two maps of the list. Then \(\beta\) for the composite is \(\beta\) for \(\Lambda'\to\Lambda''\) composed with the extension of scalars to \(\Lambda''\) of \(\beta\) for \(\Lambda\to\Lambda'\). Indeed, both are transposes of the composite of the extended counits and the comparisons for \(Rf_!\). Third, for an open immersion \(u\), under \(u^!\simeq u^*\), the map \(\beta_{u,K}\) is the isomorphism of Definition definition-change-of-coefficients-maps. Both are transposes of the same map under \(u_!\dashv u^*\).

% object: lemma-change-coefficients-exceptional-inverse-image (blueprint/chap-The-six-operations-on-Dbctf/sec-Change-of-coefficients/lemma-change-coefficients-exceptional-inverse-image/_section.md)
\phantomsection\label{bp21}%
\begin{bpstatementit}{Lemma~7.8}{ (Change of coefficients for the exceptional inverse image)}
For each map \(\Lambda\to\Lambda'\) listed above, every separated morphism \(f:X\to Y\) of finite type and every \(K\in D^b_{\mathrm{ctf}}(Y,\Lambda)\), the map \(\beta_{f,K}\) of (7.1) is an isomorphism.
\end{bpstatementit}

% proof: proof/lemma-change-coefficients-exceptional-inverse-image
\begin{proof}[{Proof}]\phantomsection\label{bp22}%
Since \(f^!\) is only known through its adjunction, there is no formula into which we could substitute \(K_{\Lambda'}\) directly. Instead we show that \(\beta_{f,K}\) is an isomorphism in two cases where \(f^!\) has been computed: closed immersions and the projection \(\mathbf P^1_Y\to Y\). We then factor a general \(f\) locally into such maps, using that \(\beta\) is compatible with composition and with restriction to opens.

\emph{Closed immersions.} For a closed immersion \(i\) with open complement \(j\), the functor \(i^!\) is determined by the localization triangle

\[
i_*i^!K\longrightarrow K\longrightarrow Rj_*j^*K.
\]

It expresses \(i_*i^!K\) as the fibre of a map between objects whose behaviour under extension of scalars we already know. Extension of scalars commutes with \(i_*\), \(j^*\) and, by Lemma \hyperref[bp16]{7.6}, \(Rj_*\). Hence the extension of scalars of this triangle is the triangle for \(K_{\Lambda'}\) over \(\Lambda'\). Since the first map is the counit, the comparison on fibres is \(i_*\beta_{i,K}\). The comparisons on the other two terms are isomorphisms, and therefore \(i_*\beta_{i,K}\) is one as well. Since \(i_*\) is fully faithful, \(\beta_{i,K}\) is an isomorphism.

\emph{The projective line.} Let \(q:\mathbf P^1_Y\to Y\). Recall that the adjunction \(Rq_*\dashv q^*(1)[2]\) between \(D^b_{\mathrm{ctf}}(\mathbf P^1_Y,\Lambda)\) and \(D^b_{\mathrm{ctf}}(Y,\Lambda)\) is defined by the diagonal class and the trace. Through these two maps, \(q^!\) is identified with \(q^*(1)[2]\). It therefore suffices to know that the two maps commute with all the maps of coefficient rings listed above. This holds by construction (Propositions \bpout{6.7} and \bpout{Proposition 6.22}). Indeed, the classes with finite coefficients are constructed from the Kummer sequence and derived extension of scalars, while the integral classes are their compatible limits, and the classes over fields are obtained from these by extension of scalars. Only one map in the list needs a further remark. For \(\mathcal O_E\to A\), the proof of Lemma \hyperref[bp16]{7.6} shows that \(A_X\) is perfect. Hence the derived tensor product with \(A_X\) commutes with the inverse limits defining the integral classes, and so the resulting classes are exactly the classes with finite coefficients. Consequently the identification \(q^!\simeq q^*(1)[2]\) is compatible with extension of scalars. Under this identification, finally, \(\beta_{q,K}\) becomes the isomorphism of Definition definition-change-of-coefficients-maps for \(q^*\), twisted by \((1)[2]\).

\emph{General \(f\).} The open immersion \(\mathbf A^1_Y\hookrightarrow\mathbf P^1_Y\) is a case of restriction to opens, and so \(\beta\) is an isomorphism for \(\mathbf A^1_Y\to Y\). Since the projection \(\mathbf A^d_Y\to Y\) is a composite of \(d\) projections of this kind, the compatibility of \(\beta\) with composition gives the case of affine space. In general, whether \(\beta_{f,K}\) is an isomorphism may be checked locally on \(X\). Using a finite affine cover, we may therefore assume that \(f\) factors as a closed immersion into an affine space over \(Y\) followed by the projection. Then the two cases already treated and compatibility with composition give the claim.
\end{proof}

% object: lemma-change-coefficients-internal-hom (blueprint/chap-The-six-operations-on-Dbctf/sec-Change-of-coefficients/lemma-change-coefficients-internal-hom/_section.md)
\phantomsection\label{bp23}%
\begin{bpstatementit}{Lemma~7.9}{ (Change of coefficients for internal Hom)}
For each map \(\Lambda\to\Lambda'\) listed above and all \(K,L\in D^b_{\mathrm{ctf}}(X,\Lambda)\), the map \(\cHom(K,L)_{\Lambda'}\to\cHom(K_{\Lambda'},L_{\Lambda'})\) of Definition definition-change-of-coefficients-maps is an isomorphism.
\end{bpstatementit}

% proof: proof/lemma-change-coefficients-internal-hom
\begin{proof}[{Proof}]\phantomsection\label{bp24}%
Internal Hom is a right adjoint. As for \(f^!\), there is no formal reason for it to commute with extension of scalars, and we argue as in the proof of Proposition \hyperref[bp5]{7.1}. The comparison map for \(\cHom(-,L)\) is compatible with triangles in the first argument, so we choose a finite stratification on whose strata \(K\) is locally \(\Lambda_S\otimes P\) with \(P\) perfect. By the five lemma and finite cones we may then reduce to the case \(K=a_!P\), with \(a:S\hookrightarrow X\) locally closed and \(P\) a locally constant perfect complex on \(S\). The identity

\[
\cHom_X(a_!P,L)
\simeq Ra_*(P^\vee\otimes^{\mathbf L}a^!L)
\]

commutes with extension of scalars by Lemmas \hyperref[bp16]{7.6} and \hyperref[bp21]{Lemma 7.8}, the isomorphism of Definition definition-change-of-coefficients-maps, and the comparison for the dual of a perfect complex. For the last of these, recall that \(P\) is locally a finite complex of finite projective modules, and for such modules \(\cHom(P,\Lambda)_{\Lambda'}=\cHom(P_{\Lambda'},\Lambda')\) term by term. Moreover, all of these identifications are defined by the same evaluation pairing. Hence their composite is the map of Definition definition-change-of-coefficients-maps.
\end{proof}

% text: prose-07-07
\phantomsection\label{bp25}%
The lemmas combine into the following statement, which is the form in which change of coefficients is used later.

% object: proposition-six-operations-change-coefficients (blueprint/chap-The-six-operations-on-Dbctf/sec-Change-of-coefficients/proposition-six-operations-change-coefficients/_section.md)
\phantomsection\label{bp26}%
\begin{bpstatementit}{Proposition~7.10}{ (The six operations and change of coefficients)}
Let \(\Lambda\to\Lambda'\) be one of the maps of coefficient rings listed above. Extension of scalars sends \(D^b_{\mathrm{ctf}}(X,\Lambda)\) into \(D^b_{\mathrm{ctf}}(X,\Lambda')\). For morphisms \(f:X\to Y\) of schemes of finite type over \(k\) (separated in the case of \(Rf_!\) and \(f^!\)), the change-of-coefficients maps of Definition \hyperref[bp12]{7.4} and (7.1) for \(f^*\), \(Rf_*\), \(Rf_!\), \(f^!\), \(\otimes^{\mathbf L}\) and \(\cHom\) are isomorphisms on objects of \(D^b_{\mathrm{ctf}}(-,\Lambda)\). For a composite \(\Lambda\to\Lambda'\to\Lambda''\) of two such maps, each of these isomorphisms is the composite of the isomorphisms for \(\Lambda\to\Lambda'\) and \(\Lambda'\to\Lambda''\). These isomorphisms are compatible with the composition isomorphisms \((gf)^*\simeq f^*g^*\), \(R(gf)_*\simeq Rg_*Rf_*\), \(R(gf)_!\simeq Rg_!Rf_!\) and \((gf)^!\simeq f^!g^!\) for composable morphisms \(f\) and \(g\). Under these isomorphisms, the extensions of scalars of the units and counits of the adjunctions \(f^*\dashv Rf_*\), \(Rf_!\dashv f^!\) and \(-\otimes^{\mathbf L}K\dashv\cHom(K,-)\) (the last counit being evaluation) over \(\Lambda\) are the corresponding units and counits over \(\Lambda'\).
\end{bpstatementit}

% proof: proof/proposition-six-operations-change-coefficients
\begin{proof}[{Proof}]\phantomsection\label{bp27}%
The first assertion is Lemma \hyperref[bp14]{7.5}. Next, the maps for \(f^*\) and \(\otimes^{\mathbf L}\) are isomorphisms by definition, and those for \(Rf_*\), \(Rf_!\), \(f^!\) and \(\cHom\) by Lemmas \hyperref[bp16]{7.6}, \hyperref[bp18]{Lemma 7.7}, \hyperref[bp21]{Lemma 7.8} and \hyperref[bp23]{Lemma 7.9}. As for the compatibilities, the one with composites of coefficient maps was explained before Lemma \hyperref[bp14]{7.5}, and for \(f^!\) after (7.1). The compatibility with composition of morphisms, on the other hand, holds for \(f^*\) by the coherence of the associativity isomorphisms for pullbacks. It holds for \(Rf_*\) because adjoint transposes of composites are composites. For \(Rf_!\) the composition isomorphism is built from \(R(qr)_*\simeq Rq_*Rr_*\) and the interchange of an open and a proper direct image. Since the interchange is a map out of an extension by zero, the claim follows from the case of \(Rf_*\) and restriction to the open subscheme. Finally, for \(f^!\) it was shown after (7.1).

For the units and counits, consider one of the three adjunctions \(F\dashv G\) over \(\Lambda\). We write \(F'\dashv G'\) for the same adjunction over \(\Lambda'\), with unit \(\eta'\) and counit \(\varepsilon'\), and write \(M'\) for \(M_{\Lambda'}\). Then the comparison for the left adjoint is an isomorphism \(c_F:F'(M')\simeq(FM)'\). In each case the comparison \(c_G:(GN)'\to G'(N')\) for the right adjoint was defined as the transpose of \((\varepsilon_N)'\circ c_F\). In other words \(\varepsilon'_{N'}\circ F'(c_G)=(\varepsilon_N)'\circ c_F\), which is the statement for the counit. For the unit we must show \(c_G\circ(\eta_M)'=G'(c_F)\circ\eta'_{M'}\) as maps \(M'\to G'((FM)')\). We compare transposes. The transpose of the right side is \(c_F\), whereas the transpose of the left side is \((\varepsilon_{FM})'\circ c_F\circ F'((\eta_M)')\). By naturality of \(c_F\) this equals \((\varepsilon_{FM}\circ F\eta_M)'\circ c_F\), which is \(c_F\) by a triangle identity.
\end{proof}

% text: prose-07-08
\phantomsection\label{bp28}%
Without the condition on Tor amplitude, change of coefficients would not preserve the categories. For instance, let \(X\) be a geometric point. The sheaf \(\mathbf F_\ell\) belongs to \(\Dbc(X,\mathbf Z/\ell^2)\) but not to \(D^b_{\mathrm{ctf}}(X,\mathbf Z/\ell^2)\). Its extension of scalars \(\mathbf F_\ell\otimes^{\mathbf L}_{\mathbf Z/\ell^2}\mathbf F_\ell\) is unbounded.

For coefficients in \(\overline{\mathbf Q}_\ell\), Lemma \bpout{3.23} expresses \(D^b_{\mathrm{ctf}}(X,\overline{\mathbf Q}_\ell)\) in terms of the categories \(D^b_{\mathrm{ctf}}(X,E)\) for finite \(E\). As a statement about ordinary categories, it reads as follows.

% object: proposition-dctf-x-qbar-2-colimit (blueprint/chap-The-six-operations-on-Dbctf/sec-Change-of-coefficients/proposition-dctf-x-qbar-2-colimit/_section.md)
\phantomsection\label{bp29}%
\begin{bpstatementit}{Proposition~7.11}{ ($\Dctf(X,\Qbar)$ as a $2$-colimit)}
Extension of scalars identifies

\[
D^b_{\mathrm{ctf}}(X,\overline{\mathbf Q}_\ell)\simeq\mathop{\text{2-colim}}_{E\subset\overline{\mathbf Q}_\ell}D^b_{\mathrm{ctf}}(X,E)
\]

as ordinary categories. In other words, every object of \(D^b_{\mathrm{ctf}}(X,\overline{\mathbf Q}_\ell)\) is isomorphic to \(K_{\overline{\mathbf Q}_\ell}\) for some finite \(E\) and some \(K\in D^b_{\mathrm{ctf}}(X,E)\), and for \(K,L\in D^b_{\mathrm{ctf}}(X,E)\) extension of scalars induces a bijection

\[
\varinjlim_{F\supset E}\Hom_{D^b_{\mathrm{ctf}}(X,F)}(K_F,L_F)
\xrightarrow{\sim}\Hom_{D^b_{\mathrm{ctf}}(X,\overline{\mathbf Q}_\ell)}(K_{\overline{\mathbf Q}_\ell},L_{\overline{\mathbf Q}_\ell}),
\]

the colimit running over finite extensions \(F\) of \(E\) in \(\overline{\mathbf Q}_\ell\).
\end{bpstatementit}

% proof: proof/proposition-dctf-x-qbar-2-colimit
\begin{proof}[{Proof}]\phantomsection\label{bp30}%
An equivalence with a \(2\)-colimit amounts to two statements, essential surjectivity on objects and a bijection on morphisms. Both are parts of Lemma \bpout{3.23}, and we recall how each is proved before putting them together.

The statement about objects is Lemma lemma-descent-finite-coefficient-field. Its proof first descends the lisse sheaves on the strata, together with their transition matrices. It then descends the connecting morphisms of a finite Postnikov tower, and in this way the complexes themselves descend. The bijection on morphisms is Lemma lemma-descent-finite-coefficient-field with \(r=0\), since \(D^b_{\mathrm{ctf}}\) is a full subcategory. Its proof uses the continuity of Ext out of a constructible complex (Lemma \bpout{3.22}), applied to the filtered colimit \(\QbarX=\varinjlim_F F_X\). Surjectivity says that every morphism descends to a finite stage, and injectivity says that two morphisms with the same image agree at a later finite stage. Now consider the category whose objects are the pairs \((E,K)\) with \(K\in D^b_{\mathrm{ctf}}(X,E)\) and whose Hom sets are these colimits. Together with descent of objects, the statements above say exactly that the functor from this category to \(D^b_{\mathrm{ctf}}(X,\overline{\mathbf Q}_\ell)\) is fully faithful and essentially surjective.
\end{proof}

% text: prose-07-09
\phantomsection\label{bp31}%
In practice the proposition is used as follows. To prove a statement about finitely many objects and morphisms of \(D^b_{\mathrm{ctf}}(X,\overline{\mathbf Q}_\ell)\), we may assume that they are extensions of scalars of objects and morphisms of \(D^b_{\mathrm{ctf}}(X,E)\) for a finite extension \(E\) of \(\mathbf Q_\ell\). Moreover, two morphisms over \(E\) that agree over \(\overline{\mathbf Q}_\ell\) already agree over a finite extension of \(E\).


% object: sec-Reduction-to-finite-coefficients (blueprint/chap-The-six-operations-on-Dbctf/sec-Reduction-to-finite-coefficients/_section.md), content hidden
\subsection*{\texorpdfstring{7.3\quad Reduction to finite coefficients}{7.3 Reduction to finite coefficients}}
\phantomsection\label{bp32}%
\addcontentsline{toc}{subsection}{\texorpdfstring{7.3\quad Reduction to finite coefficients}{7.3 Reduction to finite coefficients}}
% text: prose-07-10
\phantomsection\label{bp33}%
Many later theorems say that a canonical map is an isomorphism. We prove them with finite coefficients and then extend them to all admissible rings in three steps. First, for \(\mathcal O_E\) we reduce modulo \(\pi_E\) and use derived Nakayama. This step needs derived completeness, since \(E\) itself is a nonzero \(\mathcal O_E\)-module with \(E/\pi_E E=0\). Then for \(E\) we invert \(\ell\) on integral models, and for \(\overline{\mathbf Q}_\ell\) we descend to a finite extension. The following proposition makes this argument once and for all. It only asks that the functors \(F\), \(G\) and the map \(\theta\) of the statement commute with change of coefficients, and by Proposition \hyperref[bp26]{7.10} this holds for everything built from the six operations.

% object: proposition-isomorphisms-reduce-finite-coefficients (blueprint/chap-The-six-operations-on-Dbctf/sec-Reduction-to-finite-coefficients/proposition-isomorphisms-reduce-finite-coefficients/_section.md)
\phantomsection\label{bp34}%
\begin{bpstatementit}{Proposition~7.12}{ (Isomorphisms reduce to finite coefficients)}
For each admissible \(\Lambda\) let \(F_\Lambda,G_\Lambda:D^b_{\mathrm{ctf}}(X,\Lambda)\to D^b_{\mathrm{ctf}}(Y,\Lambda)\) be exact functors and \(\theta_\Lambda:F_\Lambda\to G_\Lambda\) a natural transformation of exact functors. Assume that for each coefficient map \(\Lambda\to\Lambda'\) of Proposition \hyperref[bp26]{7.10} there are isomorphisms \(F_\Lambda(K)\otimes\Lambda'\simeq F_{\Lambda'}(K\otimes\Lambda')\) and \(G_\Lambda(K)\otimes\Lambda'\simeq G_{\Lambda'}(K\otimes\Lambda')\) carrying \(\theta_\Lambda(K)\otimes\Lambda'\) to \(\theta_{\Lambda'}(K\otimes\Lambda')\). This holds, by Proposition \hyperref[bp26]{7.10}, whenever \(F\), \(G\) and \(\theta\) are built from the six operations, their units and counits and base change maps, covariantly in \(K\). If \(\theta_A(K)\) is an isomorphism for every finite ring \(A\) and every \(K\in D^b_{\mathrm{ctf}}(X,A)\), then \(\theta_\Lambda(K)\) is an isomorphism for every admissible \(\Lambda\) and every \(K\in D^b_{\mathrm{ctf}}(X,\Lambda)\).
\end{bpstatementit}

% proof: proof/proposition-isomorphisms-reduce-finite-coefficients
\begin{proof}[{Proof}]\phantomsection\label{bp35}%
The three kinds of coefficients are reached from finite ones in three different ways: \(\mathcal O_E\) by completion, \(E\) by inverting \(\ell\), and \(\overline{\mathbf Q}_\ell\) as a colimit over finite extensions. We treat them in this order, each case using the previous one.

Let \(\Lambda=\mathcal O_E\), and put \(C=\Cone(\theta(K))\); we must show \(C=0\). Since the complexes \(F(K)\) and \(G(K)\) lie in \(\Dbc(Y,\mathcal O_E)\), they are derived \(\pi_E\)-complete by Proposition \bpout{3.11}. Moreover, derived complete complexes are closed under cones, and so \(C\) is derived complete too. Derived Nakayama (Lemma \bpout{2.62}) therefore reduces \(C=0\) to \(C/\pi_E=0\). Now reduction modulo \(\pi_E\) is a change of coefficients along \(\mathcal O_E\to\mathcal O_E/\pi_E\). By the compatibility assumed in the statement, it identifies \(\theta(K)/\pi_E\) with \(\theta_{\mathcal O_E/\pi_E}(K/\pi_E)\). This is an isomorphism by the hypothesis for finite rings, applied to \(A=\mathcal O_E/\pi_E\). Thus \(C/\pi_E=0\), and so \(C=0\). In particular, the case \(E=\mathbf Q_\ell\) gives \(\Lambda=\mathbf Z_\ell\).

Let \(\Lambda=E\). Here completeness is of no help, since a complex over \(E\) is killed by reduction modulo \(\pi_E\). Instead we use integral models. Both \(F_E\) and \(G_E\) are exact and \(\theta_E\) is a transformation of exact functors. Hence the objects \(K\) for which \(\theta_E(K)\) is an isomorphism form a triangulated subcategory. By Lemma lemma-integral-models, every object of \(\Dbc(X,E)\) is built by finitely many cones and shifts from objects \(K_0[1/\ell]\) with \(K_0\in\Dbc(X,\mathcal O_E)\). It therefore suffices to treat these. Inverting \(\ell\) is the change of coefficients along \(\mathcal O_E\to E\). By hypothesis, \(\theta_E(K_0[1/\ell])\) is therefore identified with \(\theta_{\mathcal O_E}(K_0)[1/\ell]\). This is an isomorphism by the integral case.

Let \(\Lambda=\overline{\mathbf Q}_\ell\). By Proposition \hyperref[bp29]{7.11}, every \(K\) is isomorphic to \(K_E\otimes_E\overline{\mathbf Q}_\ell\) for some finite \(E\) and some \(K_E\in D^b_{\mathrm{ctf}}(X,E)\). By hypothesis, applied to \(E\to\overline{\mathbf Q}_\ell\), the map \(\theta_{\overline{\mathbf Q}_\ell}(K)\) is identified with \(\theta_E(K_E)\otimes_E\overline{\mathbf Q}_\ell\), which is an isomorphism by the case of \(E\). Again \(E=\mathbf Q_\ell\) is included.
\end{proof}

% text: prose-07-11
\phantomsection\label{bp36}%
After a rewriting, the proposition also applies to some contravariant constructions. For instance, the double-dual map \(K\to\mathbb D\mathbb DK\) is a natural transformation of covariant exact functors. What does not reduce in this way is the construction of an adjoint or of a map. For this reason the cycle class of the diagonal and the adjunction \(Rf_!\dashv f^*(d)[2d]\) for smooth \(f\) must be built as single morphisms with \(\mathcal O_E\)-coefficients (Sections \bpout{6.7} and \bpout{Section 9.5}). Likewise, a trace with values in \(\Lambda\) is not a map of complexes.



\clearpage\phantomsection\hypertarget{bp-outside}{}%
\section*{Outside this export}
\addcontentsline{toc}{section}{Outside this export}
References marked \ensuremath{{}^\ast} are to parts of the blueprint this document leaves out:
\begin{description}
\item[Corollary 2.29] Continuity of cohomology, in Ⅰ The six operations (\bpslug{corollary-{}\allowbreak{}continuity-{}\allowbreak{}cohomology})
\item[Lemma 2.62] Reduction and derived Nakayama, in Ⅰ The six operations (\bpslug{lemma-{}\allowbreak{}reduction-{}\allowbreak{}derived-{}\allowbreak{}nakayama})
\item[Definition 3.1] Constructible sheaves, in Ⅰ The six operations (\bpslug{definition-{}\allowbreak{}constructible-{}\allowbreak{}sheaves})
\item[Proposition 3.11] The integral constructibility criterion, in Ⅰ The six operations (\bpslug{proposition-{}\allowbreak{}integral-{}\allowbreak{}constructibility-{}\allowbreak{}criterion})
\item[Section 3.4] Finite Tor amplitude, in Ⅰ The six operations (\bpslug{sec-{}\allowbreak{}Finite-{}\allowbreak{}Tor-{}\allowbreak{}amplitude})
\item[Lemma 3.13] Flat models and generators, in Ⅰ The six operations (\bpslug{lemma-{}\allowbreak{}flat-{}\allowbreak{}models-{}\allowbreak{}generators})
\item[Lemma 3.22] Continuity of Ext, in Ⅰ The six operations (\bpslug{lemma-{}\allowbreak{}continuity-{}\allowbreak{}ext-{}\allowbreak{}out-{}\allowbreak{}constructible-{}\allowbreak{}complex})
\item[Lemma 3.23] Descent to a finite coefficient field, in Ⅰ The six operations (\bpslug{lemma-{}\allowbreak{}descent-{}\allowbreak{}finite-{}\allowbreak{}coefficient-{}\allowbreak{}field})
\item[Proposition 4.40] Proper direct image of constructible complexes, in Ⅰ The six operations (\bpslug{proposition-{}\allowbreak{}proper-{}\allowbreak{}direct-{}\allowbreak{}image-{}\allowbreak{}constructible-{}\allowbreak{}complexes})
\item[Theorem 5.2] The exceptional inverse image, in Ⅰ The six operations (\bpslug{theorem-{}\allowbreak{}exceptional-{}\allowbreak{}inverse-{}\allowbreak{}image})
\item[Proposition 5.15] Comparison with the étale derived category, in Ⅰ The six operations (\bpslug{proposition-{}\allowbreak{}comparison-{}\allowbreak{}etale-{}\allowbreak{}derived-{}\allowbreak{}category})
\item[Lemma 5.20] Inverse image and tensor product, in Ⅰ The six operations (\bpslug{lemma-{}\allowbreak{}inverse-{}\allowbreak{}image-{}\allowbreak{}tensor-{}\allowbreak{}product})
\item[Proposition 5.22] Proper direct image and extension by zero, in Ⅰ The six operations (\bpslug{proposition-{}\allowbreak{}proper-{}\allowbreak{}direct-{}\allowbreak{}image-{}\allowbreak{}extension-{}\allowbreak{}zero})
\item[Proposition 5.26] Construction and independence of $Rf_!$, in Ⅰ The six operations (\bpslug{proposition-{}\allowbreak{}construction-{}\allowbreak{}independence-{}\allowbreak{}rf-{}\allowbreak{}shriek})
\item[Proposition 5.28] Composition and base change for $Rf_!$, in Ⅰ The six operations (\bpslug{proposition-{}\allowbreak{}composition-{}\allowbreak{}base-{}\allowbreak{}change-{}\allowbreak{}rf-{}\allowbreak{}shriek})
\item[Proposition 6.7] The adjunction for the projective line, in Ⅰ The six operations (\bpslug{proposition-{}\allowbreak{}adjunction-{}\allowbreak{}projective-{}\allowbreak{}line})
\item[Proposition 6.19] Open direct image preserves finite Tor amplitude, in Ⅰ The six operations (\bpslug{proposition-{}\allowbreak{}open-{}\allowbreak{}direct-{}\allowbreak{}image-{}\allowbreak{}preserves-{}\allowbreak{}finite})
\item[Section 6.7] $\ell$-adic coefficients, in Ⅰ The six operations (\bpslug{sec-{}\allowbreak{}ell-{}\allowbreak{}adic-{}\allowbreak{}coefficients-{}\allowbreak{}2})
\item[Proposition 6.21] Direct image along an open immersion, in Ⅰ The six operations (\bpslug{proposition-{}\allowbreak{}direct-{}\allowbreak{}image-{}\allowbreak{}along-{}\allowbreak{}open-{}\allowbreak{}immersion})
\item[Proposition 6.22] The $\ell$-adic cycle class of the diagonal, in Ⅰ The six operations (\bpslug{proposition-{}\allowbreak{}ell-{}\allowbreak{}adic-{}\allowbreak{}cycle-{}\allowbreak{}class-{}\allowbreak{}diagonal})
\item[Example 6.25] Local cohomology at the origin of the affine line, in Ⅰ The six operations (\bpslug{example-{}\allowbreak{}closed-{}\allowbreak{}support})
\item[Section 6.8] Proof of the theorem, in Ⅰ The six operations (\bpslug{sec-{}\allowbreak{}Proof-{}\allowbreak{}of-{}\allowbreak{}the-{}\allowbreak{}theorem})
\item[Chapter 8] Smooth base change, in Ⅰ The six operations (\bpslug{chap-{}\allowbreak{}Smooth-{}\allowbreak{}base-{}\allowbreak{}change})
\item[Theorem 8.11] Smooth base change for all admissible coefficients, in Ⅰ The six operations (\bpslug{theorem-{}\allowbreak{}smooth-{}\allowbreak{}base-{}\allowbreak{}change-{}\allowbreak{}all-{}\allowbreak{}admissible})
\item[Section 9.5] $\ell$-adic coefficients, in Ⅱ Duality and the trace formula (\bpslug{sec-{}\allowbreak{}ell-{}\allowbreak{}adic-{}\allowbreak{}coefficients-{}\allowbreak{}4})
\item[Theorem 10.8] Verdier biduality, in Ⅱ Duality and the trace formula (\bpslug{theorem-{}\allowbreak{}verdier-{}\allowbreak{}biduality})
\end{description}
\clearpage\phantomsection%
\section*{References}
\addcontentsline{toc}{section}{References}
\bpbibitem{BBD82}{bp39}{A. Beilinson, J. Bernstein, and P. Deligne. \emph{Faisceaux pervers}, volume 100 of \emph{Astérisque}. 1982.}

\bpbibitem{BS15}{bp40}{B. Bhatt and P. Scholze. The pro-étale topology for schemes. \emph{Astérisque}, 369:99–201, 2015. \href{https://people.mpim-bonn.mpg.de/scholze/proetale.pdf}{Author's manuscript}.}

\bpbibitem{Con07}{bp41}{B. Conrad. Deligne's notes on Nagata compactifications. \emph{Journal of the Ramanujan Mathematical Society}, 22:205–257, 2007. \href{https://math.stanford.edu/\~conrad/papers/nagatafinal.pdf}{Author's manuscript}.}

\bpbibitem{Del74}{bp42}{P. Deligne. La conjecture de Weil. I. \emph{Publ. Math. IHÉS}, 43:273–307, 1974.}

\bpbibitem{Del80}{bp43}{P. Deligne. La conjecture de Weil. II. \emph{Publ. Math. IHÉS}, 52:137–252, 1980.}

\bpbibitem{Den22}{bp44}{C. Deninger. There is no “Weil-”cohomology theory with real coefficients for arithmetic curves. \emph{Ann. Sc. Norm. Super. Pisa Cl. Sci.}, 2022. To appear. arXiv:\href{https://arxiv.org/abs/2204.02714}{2204.02714}.}

\bpbibitem{Deu41}{bp45}{M. Deuring. Die Typen der Multiplikatorenringe elliptischer Funktionenkörper. \emph{Abh. Math. Sem. Univ. Hamburg}, 14:197–272, 1941.}

\bpbibitem{Eke90}{bp46}{T. Ekedahl. On the adic formalism. In \emph{The Grothendieck Festschrift, Vol. II}, volume 87 of \emph{Progress in Mathematics}, pages 197–218. Birkhäuser, 1990.}

\bpbibitem{Has36}{bp47}{H. Hasse. Zur Theorie der abstrakten elliptischen Funktionenkörper. I, II, III. \emph{J. Reine Angew. Math.}, 175:55–62, 69–88, 193–208, 1936.}

\bpbibitem{Hum72}{bp48}{J. E. Humphreys. \emph{Introduction to Lie Algebras and Representation Theory}, volume 9 of \emph{Graduate Texts in Mathematics}. Springer, 1972.}

\bpbibitem{Ill94}{bp49}{L. Illusie. Autour du théorème de monodromie locale. \emph{Astérisque}, 223:9–57, 1994. §4. \href{https://www.numdam.org/item/AST_1994__223__9_0/}{Original article}.}

\bpbibitem{Kat88}{bp50}{N. M. Katz. Travaux de Laumon. \emph{Astérisque}, 161–162:105–132, 1988. Séminaire Bourbaki, exp. 691.}

\bpbibitem{KW01}{bp51}{R. Kiehl and R. Weissauer. \emph{Weil conjectures, perverse sheaves and \(\ell\)-adic Fourier transform}, volume 42 of \emph{Ergebnisse der Mathematik}. Springer, 2001.}

\bpbibitem{Lau87}{bp52}{G. Laumon. Transformation de Fourier, constantes d'équations fonctionnelles et conjecture de Weil. \emph{Publ. Math. IHÉS}, 65:131–210, 1987.}

\bpbibitem{Mil86a}{bp53}{J. S. Milne. \href{https://www.jmilne.org/math/xnotes/AVs.pdf}{Abelian Varieties}. In \emph{Arithmetic Geometry}. 1986. Corrected online version.}

\bpbibitem{Mil86b}{bp54}{J. S. Milne. \href{https://www.jmilne.org/math/xnotes/JVs.pdf}{Jacobian Varieties}. In \emph{Arithmetic Geometry}. 1986. Corrected online version.}

\bpbibitem{Ran39}{bp55}{R. A. Rankin. Contributions to the theory of Ramanujan's function \(\tau(n)\) and similar arithmetical functions. II. \emph{Proc. Cambridge Philos. Soc.}, 35:357–372, 1939.}

\bpbibitem{SGA4}{bp37}{M. Artin, A. Grothendieck, and J.-L. Verdier, editors. \emph{Théorie des topos et cohomologie étale des schémas}, volumes 269, 270, 305 of \emph{Lecture Notes in Mathematics}. Springer, 1972–1973. SGA 4.}

\bpbibitem{SGA4½}{bp38}{P. Deligne. \emph{Cohomologie étale}, volume 569 of \emph{Lecture Notes in Mathematics}. Springer, 1977. SGA 4½; \emph{Théorèmes de finitude}, §3, and \emph{Rapport sur la formule des traces}, §3, and \emph{Sommes trigonométriques}, §1.}

\bpbibitem{SGA5}{bp56}{A. Grothendieck, editor. \emph{Cohomologie \(\ell\)-adique et fonctions \(L\)}, volume 589 of \emph{Lecture Notes in Mathematics}. Springer, 1977. SGA 5.}

\bpbibitem{SGA7II}{bp57}{P. Deligne and N. Katz, editors. \emph{Groupes de monodromie en géométrie algébrique}, volume 340 of \emph{Lecture Notes in Mathematics}. Springer, 1973. SGA 7 II; Exposé XIII, (2.1.7).}

\bpbibitem{Stacks}{bp58}{The Stacks Project Authors. \href{https://stacks.math.columbia.edu}{The Stacks Project}. Individual results are cited by their stable tags in the text.}

\bpbibitem{Var05}{bp59}{Y. Varshavsky. Lefschetz–Verdier trace formula and a generalization of a theorem of Fujiwara. 2005. Version 2. arXiv:\href{https://arxiv.org/abs/math/0505564}{math/0505564}. \href{https://arxiv.org/abs/math/0505564v2}{Author's preprint}.}

\bpbibitem{Wei48a}{bp60}{A. Weil. \emph{Sur les courbes algébriques et les variétés qui s'en déduisent}. Hermann, Paris, 1948.}

\bpbibitem{Wei48b}{bp61}{A. Weil. \emph{Variétés abéliennes et courbes algébriques}. Hermann, Paris, 1948.}

\bpbibitem{Wei49}{bp62}{A. Weil. Numbers of solutions of equations in finite fields. \emph{Bull. Amer. Math. Soc.}, 55:497–508, 1949.}

\end{document}
